\documentclass{article}
\usepackage[preprint]{neurips_2025}
\makeatletter
\renewcommand{\@noticestring}{Examen Parcial --- Redes Neuronales y Aprendizaje Profundo, UNI FIIS 2026-II.}
\makeatother
% El estilo NeurIPS fija el título "Abstract"; mismo formato, en español
\renewenvironment{abstract}{\vskip 0.075in\centerline{\large\bf Resumen}\vspace{0.5ex}\begin{quote}}{\par\end{quote}\vskip 1ex}

\usepackage[utf8]{inputenc}
\usepackage[T1]{fontenc}
\usepackage[spanish,es-noshorthands,es-tabla]{babel}
\usepackage{hyperref}
\usepackage{url}
\usepackage{booktabs}
\usepackage{amsfonts,amsmath}
\usepackage{nicefrac}
\usepackage{microtype}
\usepackage{xcolor}
\usepackage{graphicx}
\usepackage{subcaption}

% Marca los valores que dependen de los resultados finales de Colab
\newcommand{\RES}[1]{\textcolor{red}{#1}}
% Figura que no rompe la compilación si aún no existe el archivo
\newcommand{\fig}[2]{\IfFileExists{../runs/figures/#1}{\includegraphics[width=#2]{#1}}{\fbox{\parbox{#2}{\centering\small\RES{[falta \detokenize{#1}]}}}}}
\graphicspath{{../runs/figures/}}

\title{Mapeo de inundaciones con Sentinel-1 y una U-Net:\\ etiquetas doradas vs.\ débiles y generalización a Lima 2017}

\author{%
  Gabriel Ángel Nizama Duque\qquad César Luis Jiménez Miranda\\
  \textbf{Juan José Pino Reyes}\qquad \textbf{Martín Alexis Valdiviezo Palma}\\[4pt]
  Maestría en Inteligencia Artificial, Universidad Nacional de Ingeniería (FIIS)\\
  \texttt{https://github.com/gabo2607/sen1floods11-unet}
}

\begin{document}
\maketitle

\begin{abstract}
El radar de apertura sintética (SAR) de Sentinel-1 atraviesa nubes y opera de noche, por lo que es la
fuente de facto para mapear inundaciones durante la emergencia. Implementamos desde cero en PyTorch una
U-Net que segmenta agua a partir de las bandas VV y VH, y la entrenamos en Sen1Floods11. Mediante una
ablación a igual cómputo comparamos tres formas de combinar los 252 chips etiquetados a mano con 1000
chips de etiqueta débil (umbral de Otsu), y tres funciones de pérdida (BCE, Dice y BCE+Dice). El mejor
modelo alcanza un IoU de \RES{0.xx} en el test dorado, \RES{0.xx} en Bolivia (evento retenido) y apenas
\RES{0.xx} en una región de prueba propia construida para Lima (marzo de 2017) a partir de una escena
Sentinel-1 del 17 de marzo y la delineación Copernicus EMS EMSN038. Atribuimos esa brecha a tres
causas: el tipo de inundación (lodo turbulento y zonas urbanas, que no se ven oscuros en el SAR), el
desajuste de fecha y de criterio de la referencia óptica, y la geometría de radar en el relieve andino.
\end{abstract}

% ---------------------------------------------------------------------------------------------
\section{Motivación}

Saber \emph{dónde} está el agua en las horas siguientes a una inundación es crítico para la respuesta
de emergencia, pero las nubes que acompañan a las tormentas inutilizan a los satélites ópticos justo
cuando más se los necesita. El SAR en banda C de Sentinel-1~\cite{torres2012gmes} es insensible a las
nubes y adquiere de día y de noche; el agua abierta y calma se comporta como un espejo que refleja la
señal lejos del sensor y aparece oscura, lo que hace de la segmentación de agua un problema bien
planteado para redes convolucionales. El Perú tiene un precedente concreto: durante El Niño costero de
marzo de 2017 los desbordes del Rímac y del Chillón afectaron Carapongo, Huachipa, Comas, Puente Piedra
y otros distritos de Lima, y Copernicus EMS activó el mapeo EMSN038~\cite{cems2017emsn038}.

El reto no es solo arquitectónico. Las etiquetas hechas a mano son escasas (446 chips en
Sen1Floods11~\cite{bonafilia2020sen1floods11}), mientras que las etiquetas débiles generadas por
umbralización son abundantes pero ruidosas. Además, un modelo útil en el Perú debe generalizar a una
región que no vio en el entrenamiento. Este trabajo responde empíricamente a dos preguntas: (i) cómo
combinar la supervisión dorada y la débil, y (ii) cuánto rendimiento se pierde al llevar el modelo a
Lima.

% ---------------------------------------------------------------------------------------------
\section{Trabajos relacionados}

\paragraph{Umbralización clásica.} El enfoque operativo dominante umbraliza la retrodispersión,
típicamente con Otsu~\cite{otsu1979} aplicado sobre teselas bimodales, seguido de refinamiento por
crecimiento de regiones y modelos de elevación~\cite{martinis2015fully,twele2016sentinel}. Es rápido y no
requiere datos etiquetados, pero falla cuando la superficie del agua es rugosa (viento, corriente) o
cuando hay otras superficies oscuras (sombra de radar, asfalto, arena seca).

\paragraph{Aprendizaje profundo.} La U-Net~\cite{ronneberger2015unet} y sus variantes, como
Attention U-Net~\cite{oktay2018attention}, son el estándar para segmentación densa. Sen1Floods11
introdujo chips dorados y débiles de 11 eventos y mostró que redes entrenadas con etiquetas débiles
pueden superar al umbral que las generó~\cite{bonafilia2020sen1floods11}.
Konapala et al.~\cite{konapala2021exploring} estudiaron la fusión de Sentinel-1 y Sentinel-2 sobre el
mismo conjunto, y WorldFloods~\cite{mateo2021towards} mostró la escala alcanzable con etiquetas
derivadas de productos de emergencia. La pérdida Dice~\cite{milletari2016vnet} se usa con frecuencia
para contrarrestar el desbalance entre agua y tierra.

\paragraph{Etiquetas ruidosas e inundación urbana.} El aprendizaje con etiquetas ruidosas está
revisado en~\cite{frenay2014classification}; las estrategias más simples son mezclar fuentes o
preentrenar con datos ruidosos y afinar con limpios, que son las que evaluamos. En zonas urbanas, el
agua entre edificios \emph{aumenta} la retrodispersión por doble rebote en lugar de reducirla, lo que
invierte la firma del agua~\cite{mason2010urban}. Por último, la normalización radiométrica del SAR
($\sigma^0$ frente a $\gamma^0$ aplanado por terreno~\cite{small2011flattening}) cambia los niveles de
dB, sobre todo en relieve.

% ---------------------------------------------------------------------------------------------
\section{Datos}

\paragraph{Sen1Floods11.} Chips de $512\times512$ píxeles a 10\,m con $\sigma^0$ de VV y VH en dB
(EPSG:4326). Usamos las particiones oficiales de los chips dorados: 252 de entrenamiento, 89 de
validación, 90 de test (mismos eventos que el entrenamiento) y 15 de Bolivia (evento retenido). Las
etiquetas son agua, no agua y sin datos; los píxeles sin datos se ignoran en la pérdida y en las
métricas. Los valores de dB se recortan a $[-50, 1]$ y se escalan a $[0, 1]$.

\paragraph{Etiquetas débiles.} De los 4\,384 chips con etiqueta Otsu sobre S1 elegimos al azar, con
semilla fija, 1000 chips, \textbf{excluyendo todo Bolivia y cualquier ubicación presente en validación o
test} para evitar fuga de información hacia la evaluación.

\paragraph{Región de prueba propia: Lima, marzo de 2017.} Sen1Floods11 no incluye eventos peruanos,
así que construimos la región combinando dos fuentes independientes:
\begin{itemize}
  \item \textbf{SAR}: escena Sentinel-1B IW GRDH dual-pol del \textbf{17 de marzo de 2017}, órbita
  descendente. Es la escena VV+VH más cercana a la delineación y figura entre las que lista el propio
  informe de EMSN038. La obtuvimos del catálogo RTC de Microsoft Planetary Computer, que entrega
  $\gamma^0$ aplanado por terreno; la convertimos a $\sigma^0 \approx \gamma^0\cos\theta$, con el
  ángulo de incidencia $\theta$ interpolado linealmente a lo ancho del barrido IW (29.1°--46.0°), y la
  reproyectamos a la grilla de Sen1Floods11.
  \item \textbf{Referencia}: capa \emph{Potential Maximum Flood Extent} de EMSN038, delineada sobre
  WorldView-2 del 18 y 21 de marzo~\cite{cems2017emsn038}, rasterizada sobre la misma grilla. Como en
  Sen1Floods11 también se marca el agua permanente, la clase agua es la unión de la inundación, los cauces
  y las lagunas; el mar, lo que queda fuera del AOI y los píxeles sin datos se ignoran.
\end{itemize}
Conservamos los \textbf{19 chips} con al menos 500 píxeles de inundación, unos 1.9\,km² de agua de
referencia. Esta referencia es en sí misma una \emph{supervisión débil}. El propio informe de EMSN038
reconoce que el agua ya estaba en retroceso el 21 de marzo, que por la carga de sedimentos el agua es
``extremadamente difícil'' de separar del suelo desnudo, y que en el 25\,\% del AOI no se pudo delinear
la inundación. Las cifras de Lima deben leerse con esa cautela.

% ---------------------------------------------------------------------------------------------
\section{Metodología}

\paragraph{Arquitectura.} U-Net implementada a mano (\texttt{src/models/unet.py}): bloque de entrada y
4 niveles de codificador, cada uno con dos convoluciones $3\times3$ + BatchNorm + ReLU; submuestreo por
\emph{max-pooling}; decodificador con convolución transpuesta $2\times2$, concatenación con la conexión
de salto del mismo nivel y doble convolución; y una convolución $1\times1$ final que produce un logit
por píxel. Con 32 canales base (32--512) tiene 7.8\,M de parámetros. La entrada son 2 canales (VV, VH).

\paragraph{Pérdidas.} BCE sobre logits, Dice suave~\cite{milletari2016vnet} calculada por imagen y
promediada en el lote, y su combinación $\tfrac12(\text{BCE}+\text{Dice})$, todas con máscara de
píxeles válidos.

\paragraph{Entrenamiento.} AdamW~\cite{loshchilov2019adamw} (lr $10^{-3}$, decaimiento $10^{-4}$) con
\emph{cosine annealing}, precisión mixta, lotes de 8 recortes aleatorios de $256\times256$ con volteos
horizontales y verticales, y 40 épocas de 31 pasos. \textbf{Todas las variantes usan el mismo número de
pasos de optimización}, de modo que comparan estrategias de datos y no cantidad de cómputo. La
validación se hace sobre chips completos de $512\times512$, y el punto de control se elige
\emph{exclusivamente} por IoU de validación; test, Bolivia y Lima nunca intervienen en una decisión.

\paragraph{Métricas.} Para la clase agua: IoU global (agregando TP/FP/FN de todos los píxeles), IoU
promediado por chip (métrica de los autores de Sen1Floods11), F1, precisión, recall, omisión
($1-$recall), comisión ($1-$precisión) y matriz de confusión, con umbral 0.5.

% ---------------------------------------------------------------------------------------------
\section{Experimentos}

\paragraph{Ablación (a): combinación de supervisiones} (todas con BCE+Dice):
\textbf{gold}, solo los 252 chips dorados;
\textbf{mixed}, dorados y débiles mezclados uniformemente (1252 chips);
\textbf{pre\,$\to$\,ft}, 20 épocas solo con débiles y 20 épocas de afinado con dorados (lr $3\cdot10^{-4}$).

\paragraph{Ablación (b): función de pérdida.} Con la mejor estrategia de (a), elegida por validación,
se entrena con BCE sola y con Dice sola.

\paragraph{Dos rondas de entrenamiento.} En la \textbf{ronda 1} cada experimento se entrenó una vez
(semilla 0), y con ella se eligió la estrategia de datos para la ablación (b). Las curvas mostraron
que el IoU de validación oscila mucho entre épocas consecutivas (por ejemplo, de 0.61 a 0.46 en
\emph{gold}), así que una diferencia de pocas centésimas entre variantes podía deberse al azar de la
inicialización y del muestreo. En la \textbf{ronda 2} repetimos los 5 experimentos con las semillas 1 y
2. Todos los resultados se reportan como \textbf{media $\pm$ desviación estándar sobre 3
entrenamientos}, y para las figuras cualitativas se usa el entrenamiento con mejor IoU de validación del
mejor experimento. En una GPU T4, cada entrenamiento de 40 épocas tarda unos 6 minutos.

% ---------------------------------------------------------------------------------------------
\section{Resultados}

\begin{table}[t]
  \caption{Ablación. IoU, F1, precisión (P) y recall (R) de la clase agua. Test: chips dorados de los
  eventos vistos; Bolivia: evento retenido con etiquetas doradas; Lima: región propia (S1 17-mar-2017
  + EMSN038). La selección se hace solo por IoU$_{val}$.}
  \label{tab:ablation}
  \centering
  \scriptsize
  \setlength{\tabcolsep}{3.5pt}
  \IfFileExists{../runs/ablation.tex}{\input{../runs/ablation.tex}}{\RES{[falta runs/ablation.tex]}}
\end{table}

\paragraph{Ablación (a).} \RES{[Completar: cuál estrategia gana y por cuánto. Si gold gana: con
igual cómputo, el ruido de Otsu (sobre todo en sombras y suelos oscuros) cuesta más de lo que aporta la
diversidad extra. Si mixed o pre$\to$ft gana: la diversidad de escenas compensa el ruido, y afinar con
dorados corrige el sesgo del umbral.]}

\paragraph{Ablación (b).} \RES{[Completar: BCE sola tiende a mayor precisión y menor recall en clases
minoritarias; Dice sola optimiza directamente el solapamiento y sube el recall, pero puede ser
inestable en chips casi sin agua; la combinación suele equilibrar ambas.]}

\begin{figure}[t]
  \centering
  \fig{curves.png}{0.62\linewidth}
  \caption{IoU de validación por época. La línea punteada marca el paso de preentrenamiento a afinado.}
  \label{fig:curves}
\end{figure}

\paragraph{Evaluación en distribución.} El mejor modelo (\RES{nombre}) obtiene en test un IoU de
\RES{0.xx}, F1 \RES{0.xx}, precisión \RES{0.xx} y recall \RES{0.xx}. La matriz de confusión
(Fig.~\ref{fig:confusion}) muestra que el error dominante es la \textbf{omisión}: \RES{24\,\%} de los
píxeles de agua se pierden, frente a \RES{2.4\,\%} de falsos positivos sobre tierra. Como la tierra es
unas 7 veces más abundante, esos pocos falsos positivos aún pesan en la precisión. El IoU promediado por
chip (\RES{0.xx}) es bastante menor que el global, porque los chips con muy poca agua, donde cualquier
error domina el cociente, pesan igual que los chips mayormente inundados. Por evento
(Fig.~\ref{fig:events}), el rendimiento va de \RES{0.87} en Sri Lanka a \RES{0.18} en Pakistán, donde
las inundaciones son someras y con vegetación que asoma sobre el agua.

\begin{figure}[t]
  \centering
  \begin{subfigure}{0.32\linewidth}\fig{confusion_test.png}{\linewidth}\caption{Test}\end{subfigure}
  \begin{subfigure}{0.32\linewidth}\fig{confusion_bolivia.png}{\linewidth}\caption{Bolivia}\end{subfigure}
  \begin{subfigure}{0.32\linewidth}\fig{confusion_lima.png}{\linewidth}\caption{Lima}\end{subfigure}
  \caption{Matrices de confusión del mejor modelo, normalizadas por fila (etiqueta).}
  \label{fig:confusion}
\end{figure}

\begin{figure}[t]
  \centering
  \fig{events.png}{0.62\linewidth}
  \caption{IoU global por evento (test) y en las dos regiones nunca vistas.}
  \label{fig:events}
\end{figure}

% ---------------------------------------------------------------------------------------------
\section{Generalización geográfica: Bolivia y Lima}

\paragraph{Bolivia: sin brecha apreciable.} En el evento retenido el mejor modelo obtiene un IoU de
\RES{0.xx}, igual o superior al del test en distribución (\RES{0.xx}). Bolivia es una inundación fluvial
en llanura, con agua abierta y calma, del mismo tipo que la mayoría de los eventos de entrenamiento, y
sus etiquetas siguen el mismo protocolo. Por lo tanto, cambiar de región \emph{por sí solo} no degrada
al modelo cuando el fenómeno físico es el mismo.

\paragraph{Lima: colapso.} En la región peruana todos los modelos caen a un IoU de \RES{0.00x--0.0xx},
con precisión y recall por debajo de \RES{5\,\%}. Para separar el fallo del modelo de un problema de los
datos, comparamos la distribución de $\sigma^0$ de los píxeles que la referencia marca como agua frente
a los de tierra (Fig.~\ref{fig:backscatter}). En Sen1Floods11 y en Bolivia el agua es unos
\textbf{8\,dB más oscura} que la tierra en VV (medianas de $-17.6$ frente a $-9.5$\,dB en test). En Lima
las dos distribuciones \textbf{se superponen casi por completo} ($-8.7$ frente a $-9.0$\,dB en VV y
$-16.4$ frente a $-16.1$\,dB en VH). Sobre la escena del 17 de marzo, la ``inundación'' de EMSN038 no
tiene la firma del agua abierta, así que ningún clasificador que use solo SAR podría separarla de la
tierra. La brecha no se explica, entonces, por falta de capacidad del modelo, sino por las siguientes
causas:

\begin{enumerate}
  \item \textbf{Tipo de inundación y cobertura urbana.} En Lima no hubo una lámina de agua calma, sino
  huaicos y desbordes de agua turbulenta con mucho sedimento en cauces encañonados y zonas urbanas. El
  agua rugosa devuelve energía al sensor y no se ve oscura, y en la zona urbana inundada (Huachipa,
  Carapongo) el doble rebote entre la fachada y el agua \emph{aumenta} la
  retrodispersión~\cite{mason2010urban}. La Fig.~\ref{fig:backscatter} lo confirma: entre los píxeles de
  ``agua'' de Lima hay \emph{más} retornos saturados (pico cerca de $+1$\,dB) que entre los de tierra.
  Ninguna de esas firmas aparece en los eventos de entrenamiento, que son mayormente rurales y llanos.
  \item \textbf{Desajuste temporal y de criterio de la referencia.} La SAR es del 17 de marzo y la
  delineación óptica del 18 y 21, con el agua ya en retroceso. Además, la capa \emph{Potential Maximum
  Flood Extent} delinea la extensión \emph{máxima} afectada, incluidas zonas con lodo y huellas de la
  crecida, no solo el agua visible en un instante. Gran parte de la ``omisión'' es, en realidad,
  desacuerdo de criterio entre la etiqueta y lo que el sensor puede ver.
  \item \textbf{Geometría y normalización del radar (efecto menor).} El relieve de las quebradas produce
  sombra de radar e inversión por relieve (\emph{layover}), fuente de falsos positivos. En cambio, la
  conversión aproximada $\gamma^0\to\sigma^0$~\cite{small2011flattening} no parece la causa principal:
  la tierra de Lima ($-9.0$\,dB en VV) tiene niveles muy parecidos a los de Sen1Floods11 ($-9.5$\,dB).
  \item \textbf{Escala.} El agua de referencia ocupa en promedio un 3\,\% de cada chip y se reparte en
  franjas de pocas decenas de metros, al límite de la resolución de 10\,m, donde los bordes pesan mucho
  en el IoU.
\end{enumerate}

\begin{figure}[t]
  \centering
  \fig{backscatter.png}{0.95\linewidth}
  \caption{Distribución de $\sigma^0$ (dB) de los píxeles etiquetados como agua (azul) y como tierra
  (naranja). En Sen1Floods11 y en Bolivia el agua es claramente más oscura; en Lima, la referencia
  EMSN038 y la tierra son indistinguibles en ambas polarizaciones.}
  \label{fig:backscatter}
\end{figure}

% ---------------------------------------------------------------------------------------------
\section{Análisis cualitativo de errores característicos del SAR}

La Fig.~\ref{fig:errors} muestra los chips con más píxeles mal clasificados. Identificamos al menos
tres tipos de error:

\begin{enumerate}
  \item \textbf{Vegetación inundada y agua somera (omisión).} En Pakistán e India, campos anegados con
  vegetación que asoma producen dispersión de volumen y doble rebote tallo--agua, y el píxel no se ve
  oscuro. El modelo omite grandes áreas (IoU \RES{0.17} y \RES{0.28} en los peores chips).
  \item \textbf{Suelo húmedo, arena y superficies lisas (comisión).} Parcelas recién aradas o
  inundadas antes del evento, bancos de arena y suelos saturados son especularmente lisos y se ven tan
  oscuros como el agua. El modelo los marca como agua, como en el chip \RES{Pakistan\_70625}, donde la
  etiqueta solo contiene un canal estrecho.
  \item \textbf{Zonas urbanas: doble rebote.} En el chip urbano de España y en Huachipa (Lima), los
  edificios generan retornos muy brillantes que interrumpen la máscara de agua y fragmentan la
  predicción en el tejido urbano.
  \item \textbf{Sombra de radar en relieve (comisión, Lima).} En las laderas opuestas al sensor no
  llega señal, y quedan franjas oscuras que el modelo confunde con agua (Fig.~\ref{fig:errors-lima}).
\end{enumerate}

\begin{figure}[t]
  \centering
  \fig{errors_test.png}{0.92\linewidth}
  \caption{Los chips de test con más errores. Columnas: VV, VH (dB), etiqueta dorada, predicción y
  mapa de errores (gris: acierto; naranja: falso positivo; azul: falso negativo).}
  \label{fig:errors}
\end{figure}

\begin{figure}[t]
  \centering
  \fig{errors_lima.png}{0.92\linewidth}
  \caption{Chips de Lima con más errores, con la referencia EMSN038 como etiqueta.}
  \label{fig:errors-lima}
\end{figure}

% ---------------------------------------------------------------------------------------------
\section{Discusión y limitaciones}

\begin{itemize}
  \item \textbf{Presupuesto de cómputo.} Cada variante se entrenó una sola vez (semilla 0) y por 40
  épocas cortas. Diferencias de unas pocas centésimas de IoU entre variantes pueden estar dentro de la
  varianza entre semillas.
  \item \textbf{Subconjunto de etiquetas débiles.} Usamos 1000 de los 4\,384 chips débiles de S1 y no
  probamos los derivados de Sentinel-2. Con más datos débiles, o con ponderación por confianza, la
  conclusión de la ablación (a) podría cambiar.
  \item \textbf{Referencia de Lima.} Es óptica, de otra fecha y con otro criterio (extensión máxima),
  así que la brecha de Lima mezcla el error del modelo con el error de la etiqueta. Lo reportamos como
  un límite superior de la caída real.
  \item \textbf{Normalización radiométrica.} La conversión $\gamma^0\to\sigma^0$ con incidencia
  interpolada es aproximada. Lo ideal sería procesar la GRD con la misma cadena que Sen1Floods11
  (Google Earth Engine) o usar $\gamma^0$ también en el entrenamiento.
  \item \textbf{Solo SAR.} No usamos Sentinel-2 (bono opcional) ni un modelo de elevación (HAND), que
  ayudarían a descartar sombra de radar y zonas altas.
\end{itemize}

\section{Conclusiones}

Una U-Net de 7.8\,M de parámetros entrenada solo con SAR alcanza un IoU de \RES{0.xx} en distribución.
En la ablación, \RES{[estrategia]} fue la mejor forma de usar las etiquetas débiles, y \RES{[pérdida]}
la mejor pérdida. El modelo generaliza razonablemente a otra inundación fluvial y rural (Bolivia,
\RES{0.xx}), pero no al escenario peruano de 2017, cuyas firmas SAR (agua turbulenta, zonas urbanas,
relieve andino) no están en el entrenamiento. Para uso operativo en el Perú, la conclusión práctica es
que se necesitan ejemplos locales, ya sea afinando con unos pocos chips etiquetados de eventos peruanos,
y fuentes auxiliares (DEM, óptico) antes de confiar en un modelo entrenado en otras regiones.

\paragraph{Reproducibilidad.} Todo el pipeline (descarga, ablación, evaluación y figuras) corre con
\texttt{python run\_all.py}; ver el README del repositorio.

{\small
\bibliographystyle{plain}
\bibliography{refs}
}

\end{document}
